\section{训练 RNN：反向传播与梯度问题}

\subsection{RNN 的损失函数}

与前馈网络类似，RNN 的训练也需要定义损失函数。不同之处在于，RNN 在\textbf{每个时间步}都可以计算一个损失值：

$$\mathcal{L}_t = \text{Loss}(\hat{y}_t, y_t)$$

总损失是所有时间步损失的总和：

$$\mathcal{L} = \sum_{t=1}^{T} \mathcal{L}_t$$

\subsection{Backpropagation Through Time (BPTT)}

训练 RNN 使用的是反向传播的扩展版本——\textbf{时间反向传播}（Backpropagation Through Time, BPTT）。其核心思想是：

\begin{enumerate}
    \item 前向传播：逐时间步计算隐藏状态和输出
    \item 在每个时间步计算损失
    \item 反向传播：将误差沿时间轴从后向前传递，计算每个参数的梯度
    \item 更新参数
\end{enumerate}

\begin{figure}[H]
\centering
\includegraphics[width=0.85\textwidth]{slides-images/slide-45.jpg}
\caption{标准 RNN 的梯度流动\protect\footnotemark}
\end{figure}
\footnotetext{Slides 第45页。}

\subsection{梯度消失与梯度爆炸}

在 BPTT 过程中，梯度需要经过多次矩阵乘法（$W_{hh}$ 的多次连乘）和激活函数导数的连乘。这会导致两个严重的数值问题：

\begin{importantbox}{梯度消失与梯度爆炸}
\begin{itemize}
    \item \textbf{梯度爆炸}（Exploding Gradients）：当 $W_{hh}$ 的特征值 $> 1$ 时，梯度随时间步呈指数增长，导致参数更新过大，训练不稳定
    \item \textbf{梯度消失}（Vanishing Gradients）：当 $W_{hh}$ 的特征值 $< 1$ 时，梯度随时间步呈指数衰减趋近于零，使得网络无法学习长期依赖关系
\end{itemize}
\end{importantbox}

\begin{figure}[H]
\centering
\includegraphics[width=0.85\textwidth]{slides-images/slide-48.jpg}
\caption{梯度消失问题：计算关于 $h_0$ 的梯度涉及 $W_{hh}$ 的多次连乘\protect\footnotemark}
\end{figure}
\footnotetext{Slides 第48页。}

\subsection{长期依赖的困难}

梯度消失问题直接导致 RNN 难以捕捉长期依赖：

\begin{figure}[H]
\centering
\includegraphics[width=0.85\textwidth]{slides-images/slide-52.jpg}
\caption{梯度消失导致 RNN 偏向捕捉短期依赖\protect\footnotemark}
\end{figure}
\footnotetext{Slides 第52页。}

\begin{warningbox}{梯度消失 $\neq$ 梯度为零}
梯度消失并不意味着梯度完全消失，而是指来自\textbf{较远时间步}的梯度信号变得极其微弱，导致模型的参数更新几乎完全由近期时间步的梯度主导。因此，模型会偏向于学习短期依赖关系，而难以捕捉长期模式。
\end{warningbox}

\subsection{解决梯度问题的策略}

解决梯度消失/爆炸问题有三大类方法：

\begin{enumerate}
    \item \textbf{激活函数选择}：使用 ReLU 等不会在饱和区压缩梯度的激活函数
    \item \textbf{权重初始化}：精心设计权重初始化策略，避免梯度在传播过程中快速缩放
    \item \textbf{网络架构改进}：设计专门的门控机制来控制信息流动，如 LSTM 和 GRU
\end{enumerate}

\subsection{LSTM：长短期记忆网络}

\textbf{Long Short-Term Memory (LSTM)} 是解决梯度消失问题最著名的 RNN 变体。其核心思想是在标准 RNN 单元内部添加\textbf{门控机制}（Gates），用以选择性地控制信息的遗忘、存储和输出。

\begin{importantbox}{LSTM 的三个门}
\begin{itemize}
    \item \textbf{遗忘门}（Forget Gate）：决定哪些旧信息需要被丢弃
    \item \textbf{输入门}（Input Gate）：决定哪些新信息需要被存储
    \item \textbf{输出门}（Output Gate）：决定哪些信息需要被输出到下一个时间步
\end{itemize}
这些门的值由 sigmoid 函数控制，输出在 $[0, 1]$ 之间，分别表示``完全遗忘''到``完全保留''。
\end{importantbox}

LSTM 通过这种门控机制，使梯度能够沿着一条``高速公路''（Cell State）在时间步之间传递，大大缓解了梯度消失问题，使网络能够学习更长期的依赖关系。

\begin{knowledgebox}{LSTM 的历史意义}
LSTM 由 Hochreiter 和 Schmidhuber 于 1997 年提出，是深度学习历史上最重要的架构创新之一。在 Transformer 出现之前（2017年），LSTM 及其变体（如 GRU）是处理序列数据的主流选择，广泛应用于机器翻译、语音识别、文本生成等领域。
\end{knowledgebox}

\subsection{本章小结}

训练 RNN 的核心挑战是梯度消失/爆炸问题，它限制了标准 RNN 学习长期依赖的能力。LSTM 通过门控机制有效缓解了这一问题，但其本质仍然是逐时间步处理序列，存在计算效率和长序列建模能力上的局限。

\newpage
